\section{Übung 5 – Zustandsraummodell eines Masse-Feder-Systems}

Betrachtet wird ein gedämpfter Oszillator (Masse-Feder-System) mit der Masse
$M$, der Federkonstante $k_F$ und der Dämpfungskonstante $b$. Die Auslenkung
$x(t)$ und die Geschwindigkeit $v(t) = \dot{x}(t)$ bilden den Zustandsvektor
$\vec{z}(t) = (x(t),\, v(t))^{\mathsf{T}}$. Mit den Hilfsgrößen $\omega_0^2 =
\tfrac{k_F}{m}$ und $2\gamma = \tfrac{b}{m}$ lautet die Bewegungsgleichung des
freien Oszillators

$$
\frac{\mathrm{d}v(t)}{\mathrm{d}t} = -\omega_0^2\, x(t) - 2\gamma\, v(t).
\qquad (3)
$$

Wird das System zusätzlich über einen Exzenter mit der Kraft
$F(t) = \hat{F}\cos(\omega t)$ angeregt, ergibt sich

$$
\frac{\mathrm{d}v(t)}{\mathrm{d}t} = \frac{\hat{F}}{m}\cos(\omega t)
- \omega_0^2\, x(t) - 2\gamma\, v(t).
\qquad (4)
$$

% ================================================================
\subsection*{Vorüberlegung – Warum überhaupt ein Zustandsraummodell?}
% ================================================================

\begin{gedankenexperiment}[Die Grundidee]

Die Bewegungsgleichung des Oszillators ist eine \emph{Differentialgleichung
zweiter Ordnung} in $x(t)$:

$$
m\,\ddot{x}(t) = -k_F\, x(t) - b\,\dot{x}(t).
$$

Mit einer DGL zweiter Ordnung kann man schlecht „rechnen``: Sie enthält eine
zweite Ableitung, ist nicht direkt in Matrixform darstellbar und nicht ohne
Weiteres numerisch (Schritt für Schritt) lösbar.

Die zentrale Idee des Zustandsraummodells ist daher: \textbf{Wir ersetzen eine
DGL der Ordnung $n$ durch ein System aus $n$ gekoppelten DGLen erster Ordnung.}
DGLen erster Ordnung haben die Form „Ableitung $=$ Funktion der aktuellen
Werte`` und lassen sich kompakt als

$$
\frac{\mathrm{d}\vec{z}(t)}{\mathrm{d}t} = \mathbf{A}\,\vec{z}(t) + \mathbf{B}\,u(t)
$$

schreiben. Der \emph{Zustandsvektor} $\vec{z}(t)$ fasst dabei genau die
Information zusammen, die man zu einem Zeitpunkt $t$ kennen muss, um die gesamte
Zukunft des Systems zu berechnen.

\end{gedankenexperiment}

% ================================================================
\subsection*{Das universelle Rezept: Von der DGL zum Zustandsraummodell}
% ================================================================

Das folgende Vorgehen funktioniert für \emph{jede} lineare DGL mit konstanten
Koeffizienten und ist das Werkzeug für alle Aufgaben dieses Typs. Ziel ist immer
die Standardform (vgl.\ Abschnitt~\ref{sec:zustandsraum})

$$
\frac{\mathrm{d}\vec{z}(t)}{\mathrm{d}t} = \mathbf{A}\,\vec{z}(t) + \mathbf{B}\,u(t),
\qquad
\vec{y}(t) = \mathbf{C}\,\vec{z}(t) + \mathbf{D}\,u(t).
$$

\begin{bemerkung}[Kochrezept in 6 Schritten]

\textbf{Schritt 1 – Ordnung $n$ bestimmen.}
Suche die höchste vorkommende Ableitung der gesuchten Größe. Ihre Ordnung $n$
ist die Anzahl der benötigten Zustandsvariablen. \emph{Warum:} Der Zustand muss
die Funktion und alle Ableitungen bis zur $(n-1)$-ten enthalten, denn genau
diese Werte braucht man als „Startbedingungen``, um die Lösung eindeutig
fortzuschreiben.

\textbf{Schritt 2 – Zustandsvariablen wählen.}
Setze $z_1 = $ (die Größe selbst), $z_2 = \dot{z}_1$, $z_3 = \ddot{z}_1$, \ldots,
$z_n = z_1^{(n-1)}$. \emph{Warum:} Diese „Ableitungskette`` ist die natürliche
Wahl (Phasenraum) und führt direkt auf die Schritte 3 und 4.

\textbf{Schritt 3 – Die trivialen Zustandsgleichungen.}
Aus der Definition folgt unmittelbar $\dot{z}_1 = z_2$, $\dot{z}_2 = z_3$,
\ldots, $\dot{z}_{n-1} = z_n$. \emph{Warum:} Das sind keine neuen Annahmen,
sondern nur die Umbenennung „die Ableitung von $z_i$ heißt $z_{i+1}$``.

\textbf{Schritt 4 – Die „echte`` Zustandsgleichung.}
Löse die ursprüngliche DGL nach der höchsten Ableitung $z_1^{(n)} = \dot{z}_n$
auf und ersetze rechts alle Ableitungen durch die Zustandsvariablen $z_i$ (und
ggf.\ das Eingangssignal $u$). \emph{Warum:} Damit ist auch die letzte fehlende
Ableitung $\dot{z}_n$ allein durch aktuelle Werte ausgedrückt -- das System ist
geschlossen.

\textbf{Schritt 5 – Matrixform: $\mathbf{A}$ und $\mathbf{B}$ ablesen.}
Schreibe die $n$ Gleichungen $\dot{z}_i = \ldots$ untereinander und sortiere
nach den $z_j$ und nach $u$. Die Koeffizienten vor den $z_j$ bilden zeilenweise
die Systemmatrix $\mathbf{A}$, die Koeffizienten vor $u$ die Eingangsmatrix
$\mathbf{B}$. \emph{Warum:} Das ist nur das Umschreiben eines linearen
Gleichungssystems in Matrix-Vektor-Schreibweise.

\textbf{Schritt 6 – Ausgangsgleichung: $\mathbf{C}$ und $\mathbf{D}$.}
Frage: \emph{Welche Größe wird gemessen/beobachtet?} Drücke diese durch die
Zustände (und ggf.\ $u$) aus. Die Koeffizienten bilden $\mathbf{C}$ (vor $z_j$)
und $\mathbf{D}$ (vor $u$). \emph{Warum:} Der Zustand $\vec{z}$ ist intern; das
Modell muss noch sagen, was davon nach außen sichtbar ist.

\textbf{Kontrolle:} Dimensionen müssen passen --
$\mathbf{A}\!:\,n\times n$, $\mathbf{B}\!:\,n\times L$,
$\mathbf{C}\!:\,M\times n$, $\mathbf{D}\!:\,M\times L$
($n$ Zustände, $L$ Eingänge, $M$ Ausgänge).

\end{bemerkung}

% ================================================================
\subsection*{Teil a) – Zustandsraummodell des freien Oszillators}
% ================================================================

\begin{beispiel}[Zustandsraummodell aus DGL (3)]

\textbf{Gegeben:} Die DGL (3) des freien gedämpften Oszillators,
$\dot{v}(t) = -\omega_0^2\, x(t) - 2\gamma\, v(t)$, mit $v(t) = \dot{x}(t)$.
Vorgegebener Zustandsvektor $\vec{z}(t) = (x(t),\, v(t))^{\mathsf{T}}$. Gemessen
wird die Auslenkung $x(t)$ (Skala in Bild~1).

\textbf{Gesucht:} Das vollständige Zustandsraummodell, d.\,h.\ die Matrizen
$\mathbf{A}$, $\mathbf{B}$, $\mathbf{C}$, $\mathbf{D}$ der Zustands- und
Ausgangsgleichung.

\textbf{Formel} (zeitinvariantes Zustandsraummodell,
Abschnitt~\ref{sec:zustandsraum}):

$$
\frac{\mathrm{d}\vec{z}(t)}{\mathrm{d}t} = \mathbf{A}\,\vec{z}(t) + \mathbf{B}\,u(t),
\qquad
y(t) = \mathbf{C}\,\vec{z}(t) + \mathbf{D}\,u(t).
$$

\textbf{Rechnung:} Wir gehen das Kochrezept Schritt für Schritt durch.

\medskip
\textit{Schritt 1 – Ordnung.}
Die zugrunde liegende Bewegungsgleichung $m\ddot{x} = -k_F x - b\dot{x}$ enthält
die zweite Ableitung $\ddot{x}$. Die Ordnung ist also $n = 2$
$\Rightarrow$ wir brauchen \textbf{zwei} Zustandsvariablen. Dies ist in der
Aufgabe bereits durch $\vec{z} = (x, v)^{\mathsf{T}}$ vorweggenommen.

\medskip
\textit{Schritt 2 – Zustandsvariablen.}
Gemäß der Ableitungskette wählen wir

$$
z_1(t) = x(t) \quad\text{(Auslenkung)}, \qquad
z_2(t) = \dot{x}(t) = v(t) \quad\text{(Geschwindigkeit)}.
$$

\emph{Woher:} Position und Geschwindigkeit sind genau die zwei Größen, die man
zu einem Zeitpunkt kennen muss, um die weitere Schwingung eindeutig zu bestimmen
(physikalischer Anfangszustand).

\medskip
\textit{Schritt 3 – Triviale Zustandsgleichung.}
Aus $z_2 = \dot{z}_1$ folgt unmittelbar

$$
\dot{z}_1(t) = \dot{x}(t) = v(t) = z_2(t).
$$

\emph{Warum:} Das ist keine Physik, sondern reine Definition -- die erste Zeile
des Modells „verkettet`` nur $x$ und $v$.

\medskip
\textit{Schritt 4 – Echte Zustandsgleichung.}
Die zweite Zeile ist die gegebene DGL (3), bereits nach der höchsten Ableitung
$\dot{v} = \ddot{x}$ aufgelöst:

$$
\dot{z}_2(t) = \dot{v}(t) = -\omega_0^2\, x(t) - 2\gamma\, v(t)
            = -\omega_0^2\, z_1(t) - 2\gamma\, z_2(t).
$$

\emph{Woher:} direkt aus der Bewegungsgleichung; rechts wurde nur
$x \to z_1$, $v \to z_2$ eingesetzt.

\medskip
\textit{Schritt 5 – Matrixform.}
Wir schreiben beide Zeilen untereinander und sortieren nach $z_1, z_2$:

$$
\begin{aligned}
\dot{z}_1 &= 0\cdot z_1 + 1\cdot z_2, \\
\dot{z}_2 &= -\omega_0^2\cdot z_1 - 2\gamma\cdot z_2.
\end{aligned}
$$

Die Koeffizienten in die Matrix übertragen (erste Zeile: $0,\,1$; zweite Zeile:
$-\omega_0^2,\,-2\gamma$):

$$
\frac{\mathrm{d}}{\mathrm{d}t}
\begin{pmatrix} z_1 \\ z_2 \end{pmatrix}
=
\begin{pmatrix} 0 & 1 \\ -\omega_0^2 & -2\gamma \end{pmatrix}
\begin{pmatrix} z_1 \\ z_2 \end{pmatrix}.
$$

Der freie Oszillator hat \emph{kein} Eingangssignal, daher $\mathbf{B} = \vec{0}$
und der Term $\mathbf{B}\,u(t)$ entfällt.

\medskip
\textit{Schritt 6 – Ausgangsgleichung.}
Gemessen wird laut Aufgabe nur die Auslenkung, also $y(t) = x(t) = z_1(t)$:

$$
y(t) = \underbrace{(1\;\; 0)}_{\mathbf{C}}
\begin{pmatrix} z_1 \\ z_2 \end{pmatrix} = z_1(t),
\qquad \mathbf{D} = 0.
$$

\emph{Warum:} Der Zustand enthält auch die (nicht abgelesene) Geschwindigkeit
$z_2$; $\mathbf{C}$ wählt aus, dass davon nur $z_1 = x$ sichtbar ist.

\textbf{Lösung:}

$$
\boxed{
\frac{\mathrm{d}\vec{z}(t)}{\mathrm{d}t}
= \underbrace{\begin{pmatrix} 0 & 1 \\ -\omega_0^2 & -2\gamma \end{pmatrix}}_{\mathbf{A}}
\,\vec{z}(t),
\qquad
y(t) = \underbrace{(1\;\; 0)}_{\mathbf{C}}\,\vec{z}(t)
}
$$

\textit{Kontrolle:} $\mathbf{A}$ ist $2\times 2$, $\mathbf{C}$ ist $1\times 2$ --
passend zu $n=2$ Zuständen und $M=1$ Ausgang. \checkmark

Das System ist \emph{autonom} (homogen): Ohne äußere Anregung wird die Dynamik
allein durch $\mathbf{A}$ und den Anfangszustand $\vec{z}_0$ bestimmt. Die
abklingende Schwingung in Bild~1 ist die Eigenbewegung aus einer impulsartigen
Anfangsauslenkung.

\end{beispiel}

% ================================================================
\subsection*{Teil b) – Zustandsraummodell des angeregten Oszillators}
% ================================================================

\begin{beispiel}[Zustandsraummodell aus DGL (4)]

\textbf{Gegeben:} Die DGL (4) des fremderregten Oszillators,
$\dot{v}(t) = \tfrac{\hat{F}}{m}\cos(\omega t) - \omega_0^2 x(t) - 2\gamma v(t)$,
mit der Anregungskraft $F(t) = \hat{F}\cos(\omega t)$ als Eingangssignal.
Zustandsvektor $\vec{z}(t) = (x(t),\, v(t))^{\mathsf{T}}$.

\textbf{Gesucht:} Das vollständige Zustandsraummodell mit $\mathbf{A}$,
$\mathbf{B}$, $\mathbf{C}$, $\mathbf{D}$.

\textbf{Formel} (zeitinvariantes Zustandsraummodell,
Abschnitt~\ref{sec:zustandsraum}):

$$
\frac{\mathrm{d}\vec{z}(t)}{\mathrm{d}t} = \mathbf{A}\,\vec{z}(t) + \mathbf{B}\,u(t),
\qquad
y(t) = \mathbf{C}\,\vec{z}(t) + \mathbf{D}\,u(t).
$$

\textbf{Rechnung:} Dasselbe Rezept; neu ist nur, dass jetzt ein Eingangssignal
existiert.

\medskip
\textit{Schritt 1+2 – Ordnung und Zustandsvariablen.}
Unverändert: Die DGL ist weiterhin zweiter Ordnung in $x$, also $n = 2$ mit

$$
z_1(t) = x(t), \qquad z_2(t) = v(t).
$$

\medskip
\textit{Eingangssignal identifizieren.}
Die Aufgabe sagt ausdrücklich: „Die zusätzliche Kraft stellt das Eingangssignal
dar.`` Also setzen wir

$$
u(t) = F(t) = \hat{F}\cos(\omega t).
$$

\emph{Warum wichtig:} In Schritt 5 müssen wir die Terme nach Zuständen $z_j$
\emph{und} Eingang $u$ trennen -- daher muss $u$ vorher eindeutig benannt sein.

\medskip
\textit{Schritt 3 – Triviale Zustandsgleichung.}
Wie in Teil a):

$$
\dot{z}_1(t) = v(t) = z_2(t).
$$

\medskip
\textit{Schritt 4 – Echte Zustandsgleichung.}
DGL (4) nach $\dot{v}$ aufgelöst und der Anregungsterm über
$\tfrac{\hat{F}}{m}\cos(\omega t) = \tfrac{1}{m}F(t) = \tfrac{1}{m}u(t)$
geschrieben:

$$
\dot{z}_2(t) = \dot{v}(t)
= -\omega_0^2\, z_1(t) - 2\gamma\, z_2(t) + \frac{1}{m}\,u(t).
$$

\emph{Woher kommt der $\tfrac{1}{m}$-Term:} Aus dem zweiten Newtonschen Gesetz
$F_{\text{ges}} = m\,\dot{v}$ folgt für die \emph{äußere} Kraft der Beitrag
$\dot{v} = \ldots + \tfrac{1}{m}F(t)$ -- eine Kraft wirkt über $a = F/m$ auf die
Geschwindigkeitsänderung.

\medskip
\textit{Schritt 5 – Matrixform.}
Beide Zeilen, sortiert nach $z_1, z_2$ und $u$:

$$
\begin{aligned}
\dot{z}_1 &= 0\cdot z_1 + 1\cdot z_2 + 0\cdot u, \\
\dot{z}_2 &= -\omega_0^2\cdot z_1 - 2\gamma\cdot z_2 + \tfrac{1}{m}\cdot u.
\end{aligned}
$$

Koeffizienten vor $z_1, z_2$ $\to \mathbf{A}$; Koeffizienten vor $u$ $\to
\mathbf{B}$:

$$
\frac{\mathrm{d}}{\mathrm{d}t}
\begin{pmatrix} z_1 \\ z_2 \end{pmatrix}
=
\begin{pmatrix} 0 & 1 \\ -\omega_0^2 & -2\gamma \end{pmatrix}
\begin{pmatrix} z_1 \\ z_2 \end{pmatrix}
+
\begin{pmatrix} 0 \\ \tfrac{1}{m} \end{pmatrix}
u(t).
$$

\emph{Beachte:} Die erste Komponente von $\mathbf{B}$ ist $0$, weil die Kraft
\emph{nicht direkt} auf die Position $z_1$ wirkt, sondern erst über die
Beschleunigung auf $z_2$.

\medskip
\textit{Schritt 6 – Ausgangsgleichung.}
Gemessen wird weiterhin die Auslenkung, $y(t) = x(t) = z_1(t)$:

$$
y(t) = (1\;\; 0)\,\vec{z}(t), \qquad \mathbf{D} = 0.
$$

\textbf{Lösung:}

$$
\boxed{
\frac{\mathrm{d}\vec{z}(t)}{\mathrm{d}t}
= \underbrace{\begin{pmatrix} 0 & 1 \\ -\omega_0^2 & -2\gamma \end{pmatrix}}_{\mathbf{A}}
\,\vec{z}(t)
+ \underbrace{\begin{pmatrix} 0 \\ \tfrac{1}{m} \end{pmatrix}}_{\mathbf{B}}
F(t),
\qquad
y(t) = \underbrace{(1\;\; 0)}_{\mathbf{C}}\,\vec{z}(t)
}
$$

\textit{Kontrolle:} $\mathbf{A}\!:\,2\times2$, $\mathbf{B}\!:\,2\times1$,
$\mathbf{C}\!:\,1\times2$ -- passend zu $n=2$, $L=1$, $M=1$. \checkmark

\textbf{Vergleich mit Teil a):} Die Systemmatrix $\mathbf{A}$ ist
\emph{identisch} -- die innere Dynamik (Eigenfrequenz $\omega_0$, Dämpfung
$\gamma$) ändert sich durch die Anregung nicht. Hinzugekommen ist allein die
Eingangsmatrix $\mathbf{B}$, die beschreibt, \emph{wie} das Eingangssignal in das
System eingreift. Die erzwungene Schwingung in Bild~2 ist daher die Summe aus
abklingender Eigenschwingung (homogene Lösung, $\mathbf{A}$) und der stationären
Antwort auf die periodische Anregung (partikuläre Lösung, $\mathbf{B}\,u$).

\textbf{Bemerkung:} Wählt man statt der Kraft das normierte Signal
$u(t) = \cos(\omega t)$ als Eingang, so wandert die Amplitude in die
Eingangsmatrix: $\mathbf{B} = (0,\; \tfrac{\hat{F}}{m})^{\mathsf{T}}$. Beide
Darstellungen beschreiben dasselbe System; sie unterscheiden sich nur in der
Frage, ob der Faktor $\hat{F}$ zum Signal $u$ oder zur Matrix $\mathbf{B}$
gezählt wird.

\end{beispiel}
